\documentclass[10pt,a4paper]{amsart}
\usepackage[margin=19mm]{geometry}
\usepackage{amsmath,amssymb,mathtools}
\usepackage[scr=rsfs,cal=euler]{mathalfa}
\usepackage{xfrac}
\usepackage[unicode,hidelinks,bookmarks=false]{hyperref}
\newcommand{\ie}{i.e.\ }
\newcommand{\cf}{cf.\ }
\newcommand{\resp}{respectively\ }
\newcommand{\Subsection}{\subsection}
\providecommand{\nobreakdash}{\nobreak\hbox{-}}
\setlength{\emergencystretch}{2em}
\multlinegap0pt
\usepackage{bm}
\usepackage{bbm}
\usepackage{dsfont}
\usepackage{xspace}
\usepackage{cancel}
\usepackage{mathtools}
\usepackage{amsthm}
\makeatletter
\@ifundefined{thm}{\newtheorem{thm}{Thm}}{}
\@ifundefined{lem}{\newtheorem{lem}[thm]{Lemma}}{}
\makeatother
\title[Computing the action of a matrix exponential on an interval ]{Computing the action of a matrix exponential on an interval via the $\star$-product approach}
\author{Stefano Pozza, Shazma Zahid}
\date{Source: arXiv:2602.17516v2 (CC BY 4.0)}
\begin{document}
\maketitle

\noindent\emph{Source and licence.} Computing the action of a matrix exponential on an interval via the $\star$-product approach. Version arXiv:2602.17516v2 (CC BY 4.0), distributed under the Creative Commons Attribution 4.0 International licence, as recorded in the arXiv licence metadata for this exact version and shown on the abstract page https://arxiv.org/abs/2602.17516v2. Full rights notice: licenses/arxiv-2602.17516-CC-BY-4.0.txt.\par\smallskip

\noindent\textit{Editorial scope.} Counted unit: the Lem whose source label is \texttt{lem:bound:mtx} in the article (source0.tex lines 471--480, source label \texttt{lem:bound:mtx}), delivered together with the article's own complete original proof of it (source0.tex lines 481--524, one \texttt{proof} environment of 44 lines). No mathematics is written, repaired or re-derived: statement and proof are the article's own text line for line. Required context retained uncounted: eq:pl:star (lines 366--368); eq:pl:mtx (lines 370--372); lem:bound:star (lines 382--388); lem:tridiag (lines 435--444). Every definition, display and label of those lines is retained unchanged. Omitted from this package and not redistributed: 1--365, 369--369, 373--381, 389--434, 445--470, 525--873. Nothing inside a retained excerpt is omitted or altered: every retained body is delivered exactly as the article's lines are written, so no omission receipt of any kind accompanies this package. Citations: the retained excerpts carry no citation command at all, so no bibliography entry of the article is transcribed and no reference is redistributed. Reference adaptation: none was needed and none was made; every reference inside the delivered unit resolves inside it, so no unresolved reference or citation is produced. Printed slips kept literal: no printed word, symbol, hypothesis or number of the counted statement or of its proof was altered. Layout and class: the article's own class is \texttt{elsarticle} with packages including amssymb, booktabs, graphicx, subcaption, amsmath, xcolor, amsthm; the wrapper is the lane's pinned \texttt{amsart} editorial wrapper at 10pt on A4, which restates the theorem-like environments the retained text needs and adds the article's own macro definitions that the retained text needs (none), with extra wrapper packages (bm, bbm, dsfont, xspace, cancel, mathtools). The delivered numbering is the unit's own. Assets: no retained span contains a figure environment, a graphics-inclusion command, a PDF-page inclusion, a table or a TikZ picture; no image file is copied into this package, the package declares no asset dependency and no image file is redistributed with it. Licence: version arXiv:2602.17516v2 of this article is distributed under the Creative Commons Attribution 4.0 International licence, as recorded in the arXiv licence metadata for this exact version and shown on the abstract page https://arxiv.org/abs/2602.17516v2. A full rights notice is delivered at licenses/arxiv-2602.17516-CC-BY-4.0.txt. Characters: every retained excerpt is pure ASCII, so the delivered engine, XeLaTeX with its default Latin Modern text font, reports no missing character.
\begin{equation}\label{eq:pl:star}
 p_\ell^\star(A)(t,s) := \sum_{k=0}^{\ell} \frac{ \left( A \Theta(t-s) - c \delta(t-s) \right)^{\star k}}{(1-c)^{k+1}}. 
 \end{equation}

\begin{equation}\label{eq:pl:mtx}
 p_\ell(A \otimes T_\infty) := \sum_{k=0}^{\ell} \frac{ \left( A \otimes T_\infty - c I \right)^{k}}{(1-c)^{k+1}},
\end{equation}

\begin{lem}\label{lem:bound:star} Let $A$ be a diagonalizable matrix with spectral radius $\rho(A)$ and eigenvector basis $Z$, and let $p_\ell^{\star}$ be defined as in \eqref{eq:pl:star}. Then, for 
 \[|c| < |1-c|, \quad \ell > 4\rho(A)\cdot(t+1) \frac{1}{|c|}\left(\ln\left(\frac{|1-c|}{|c|}\right)\right)^{-2} - 1, \quad \gamma_\ell :=\exp\left(2\sqrt{\frac{|\lambda|(t+1)}{(\ell+1)|c|}}\right), \]
the following bound holds
\[
      \| \exp(A(t\hspace{-2pt}+\hspace{-2pt}1)) -  [p_\ell^\star(A)\star \Theta](t,-1) \| \leq \kappa(Z)  \frac{1}{|1-c|-|c|\gamma_\ell}  \left|\frac{c \gamma_\ell}{1-c}\right|^{\ell+1}, \]
where $\kappa(Z) = \| Z \| \, \| Z^{-1}\|$ is the condition number of the eigenvector basis.
\end{lem}

\begin{lem}\label{lem:tridiag}
Let $T_\infty$ and $\hat{T}_{M}$ be the matrices defined above. Then
\begin{equation}\label{eq:lem:matching}
     T_\infty^k\, e_1 =  \hat{T}_{M}^k\, e_1, \quad k = 0, \dots, M-1.
\end{equation}
Moreover,
\begin{equation}\label{eq:lem:norm:bound}
   \left\| T_\infty^k\, e_1 \sqrt{2} \right\| \leq \frac{28}{k!} \left( \frac{9}{4} \right)^k, \quad k=0, 1, 2, \dots \,.
\end{equation}
\end{lem}

\begin{lem}\label{lem:bound:mtx} Let $A$ be a diagonalizable matrix with spectral radius $\rho(A)$ and eigenvector basis $Z$, and let $p_\ell$ be defined as in \eqref{eq:pl:mtx}. Moreover, assume that there exists $c \in \mathbb{C}$ so that \((I-A\otimes T_{M})^{-1}\) can be expanded into a convergent Taylor series centered at $c$.
Then, for 
 \[ 9\rho(A) \left(\ln\left(\frac{|1-c|}{|c|}\right)\right)^{-2}\frac{1}{|c|} - 1 \leq \ell \leq M-2, \quad \gamma_\ell :=\exp\left(\sqrt{\frac{9\rho(A)}{(\ell+1)|c|}}\right), \]
the following bound holds
\begin{align*}
\left\|(I\hspace{-1pt}\otimes \hspace{-1pt}\phi(t)^T)\left( (I\hspace{-2pt}-\hspace{-2pt}A\hspace{-1pt}\otimes\hspace{-1pt}\hat{T}_{M})^{-1}\hspace{-4pt}-\hspace{-3pt}p_\ell(A\hspace{-1pt}\otimes\hspace{-1pt}\hat{T}_{M}) \right) (I\hspace{-1pt}\otimes\hspace{-1pt} e_1\hspace{-5pt} \sqrt{2})\right\|\hspace{-2pt} 
&\leq 28 C_M \kappa(Z) M \left| \frac{c\gamma_\ell}{1-c} \right| ^{\ell + 1}, 
\end{align*}
with \(C_M:= \| (I-A \otimes \hat{T}_{M})^{-1} \|\).
\end{lem}

\begin{proof}
The matrix \( (I - A \otimes T_{M})^{-1}  (I \otimes e_1 \sqrt{2})\) has finite size, while the matrix \( (I - A \otimes \hat{T}_{M})^{-1}  (I \otimes e_1 \sqrt{2})\) has an infinite number of rows. However, this difference is virtual, since all the additional rows of \( (I - A \otimes \hat{T}_{M})^{-1}  (I \otimes e_1 \sqrt{2})\) are equal to zero.
Therefore, the Taylor expansions of \( (I - A \otimes T_{M})^{-1}  (I \otimes e_1 \sqrt{2})\) and \( (I - A \otimes \hat{T}_{M})^{-1}  (I \otimes e_1 \sqrt{2})\) converge under the same conditions on \(c\). Then, the well-known remainder of the Taylor expansion gives
\begin{align*}
\left( (I\hspace{-2pt}-\hspace{-2pt}A\hspace{-1pt}\otimes\hspace{-1pt}\hat{T}_{M})^{-1}\hspace{-4pt}-\hspace{-3pt}p_\ell(A\hspace{-1pt}\otimes\hspace{-1pt}\hat{T}_{M}) \right) (I\hspace{-1pt}\otimes\hspace{-1pt} e_1\hspace{-5pt} \sqrt{2})\hspace{-2pt} &=  (I-A \otimes \hat{T}_{M})^{-1}   \left( \frac{A\otimes \hat{T}_{M} - cI}{1-c} \right)^{\ell+1} (I \otimes e_1 \sqrt{2}) \\
&=  (I-A \otimes \hat{T}_{M})^{-1}   \left( \frac{A\otimes {T}_{\infty} - cI}{1-c} \right)^{\ell+1} (I \otimes e_1 \sqrt{2}) \\
&=  \frac{(I\hspace{-2pt}-\hspace{-2pt}A\hspace{-1pt}\otimes\hspace{-1pt}\hat{T}_{M})^{-1}}{(1-c)^{\ell+1}} \hspace{-2pt} \sum_{k=0}^{\ell+1} \binom{\ell+1}{k} \left( A^k\otimes {T}_{\infty}^k e_1\hspace{-5pt} \sqrt{2}\right) (-c)^{\ell+1-k} 
\end{align*}
%Recall that our focus is on the truncated problem \((I \otimes  \phi(t)^T) (I-A \otimes T_{M})^{-1}(I\otimes e_1 \sqrt{2})\), where $T_{M}$ is a matrix of size \(M \times M\) and \(\phi(t), e_1\) are here vectors of size $M$. 
Moreover, denoting by \(\hat{\phi}^{(M)}(t)\) the vector obtained by setting to zero all elements after the $M$-th one in $\phi(t)$, we obtain
\begin{align*} %\label{eq:equiv:trunc1}
    (I\hspace{-1pt}\otimes \hspace{-1pt}\hat\phi^{(M)}(t)^T) (I\hspace{-2pt}-\hspace{-2pt}A\hspace{-1pt}\otimes\hspace{-1pt}\hat{T}_{M})^{-1}(I\hspace{-1pt}\otimes\hspace{-1pt} e_1\hspace{-5pt} \sqrt{2}) 
    &= (I\hspace{-1pt}\otimes \hspace{-1pt}\phi(t)^T) (I\hspace{-2pt}-\hspace{-2pt}A\hspace{-1pt}\otimes\hspace{-1pt}\hat{T}_{M})^{-1}  (I\hspace{-1pt}\otimes\hspace{-1pt} e_1\hspace{-5pt} \sqrt{2}) \\ %\label{eq:equiv:trunc2}
    (I\hspace{-1pt}\otimes \hspace{-1pt}\hat\phi^{(M)}(t)^T)p_\ell(A\hspace{-1pt}\otimes\hspace{-1pt}\hat{T}_{M}) (I\hspace{-1pt}\otimes\hspace{-1pt} e_1\hspace{-5pt} \sqrt{2}) 
    &= (I\hspace{-1pt}\otimes \hspace{-1pt}\phi(t)^T) p_\ell(A\hspace{-1pt}\otimes\hspace{-1pt}\hat{T}_{M})(I\hspace{-1pt}\otimes\hspace{-1pt} e_1\hspace{-5pt} \sqrt{2}),
\end{align*}
%Thanks to Eq.~\eqref{eq:equiv:trunc1}-\eqref{eq:equiv:trunc2},
where the second equality holds since $\ell \leq M-2$. As a consequence, 
\begin{align*}
G_\ell(t) &:= (I\hspace{-1pt}\otimes \hspace{-1pt}\phi(t)^T)\left( (I\hspace{-2pt}-\hspace{-2pt}A\hspace{-1pt}\otimes\hspace{-1pt}\hat{T}_{M})^{-1}\hspace{-4pt}-\hspace{-3pt}p_\ell(A\hspace{-1pt}\otimes\hspace{-1pt}\hat{T}_{M}) \right) (I\hspace{-1pt}\otimes\hspace{-1pt} e_1\hspace{-5pt} \sqrt{2}) \\ 
&=     (I\hspace{-1pt}\otimes \hspace{-1pt}\hat{\phi}^{(M)}(t)^T)\left( (I\hspace{-2pt}-\hspace{-2pt}A\hspace{-1pt}\otimes\hspace{-1pt}\hat{T}_{M})^{-1}\hspace{-4pt}-\hspace{-3pt}p_\ell(A\hspace{-1pt}\otimes\hspace{-1pt}\hat{T}_{M}) \right) (I\hspace{-1pt}\otimes\hspace{-1pt} e_1\hspace{-5pt} \sqrt{2}).
\end{align*}
Thus, combining the previous relations with Lemma~\ref{lem:tridiag} gives
\begin{align*}
\left\|G_\ell(t)\right\|
&\leq  \frac{\| (I \otimes \hat\phi^{(M)}(t)^T) (I-A \otimes \hat{T}_{M})^{-1} \|}{|(1-c)^{\ell+1}|}  \sum_{k=0}^{\ell+1} \binom{\ell+1}{k} \left\| A^k\otimes {T}_{\infty}^k e_1 \sqrt{2}\right\| |c|^{\ell+1-k} \\
&\leq  \frac{\| I \otimes \hat\phi^{(M)}(t)^T\| \, C_M}{|(1-c)^{\ell+1}|}  \sum_{k=0}^{\ell+1} \binom{\ell+1}{k} \left\| A^k \right\| \left\| {T}_{\infty}^k  e_1 \sqrt{2} \right\| |c|^{\ell+1-k} \\
&\leq  \frac{\|\hat\phi^{(M)}(t) \|C_M}{|(1-c)^{\ell+1}|}  \sum_{k=0}^{\ell+1} \binom{\ell+1}{k} \left\| A \right\|^k  \frac{28}{k!} \left( \frac{9}{4} \right)^k |c|^{\ell+1-k} \\
%&\leq  \frac{5 C_{M}}{|(1-c)^{\ell+1}|} \hspace{-2pt} \sum_{k=0}^{\ell+1} \binom{\ell+1}{k}   \frac{\left\| A \right\|^k}{k!} \left( \frac{9}{4} \right)^k |c|^{\ell+1-k} \\
%&\leq  \kappa(Z) \frac{5 C_{M}}{|(1-c)^{\ell+1}|} \hspace{-2pt} \sum_{k=0}^{\ell+1} \binom{\ell+1}{k}   \frac{1}{k!} \left( \frac{9 \rho(A)}{4} \right)^k |c|^{\ell+1-k} \\
&\leq 28 M C_M \kappa(Z) \frac{|c|^{\ell+1}}{|(1-c)^{\ell+1}|}  \sum_{k=0}^{\ell+1} \binom{\ell+1}{k}   \frac{1}{k!} \left( \frac{9 \rho(A)}{4 |c|} \right)^k \\
&\leq 28 \kappa(Z) C_M M \left( \left| \frac{c}{(1-c)} \right| \ \exp \left( \sqrt{\frac{9 \rho(A)}{(\ell + 1)|c|}} \right) \right)^{\ell + 1},
\end{align*}
where we used the fact that \( \| \phi^{(M)}(t) \| \leq M\) since the $j$-th element of the vector is bounded by \(|[\phi^{(M)}(t)]_j| \leq \sqrt{2j+1}\), \(t \in [-1,1]\).
Similarly to the proof of Lemma~\ref{lem:bound:star}, for
\[\ell+1 > 9\rho(A) \left(\ln\left(\frac{|1-c|}{|c|}\right)\right)^{-2}\frac{1}{|c|}, \quad\gamma_\ell :=\exp\left(\sqrt{\frac{9\rho(A)}{(\ell+1)|c|}}\right), \] 
%Therefore, for \(\ell \geq 9 \rho(A)/|c| -1 \) 
 it holds
\begin{align*}
\left\|(I\hspace{-1pt}\otimes \hspace{-1pt}\phi(t)^T)\left( (I\hspace{-2pt}-\hspace{-2pt}A\hspace{-1pt}\otimes\hspace{-1pt}\hat{T}_{M})^{-1}\hspace{-4pt}-\hspace{-3pt}p_\ell(A\hspace{-1pt}\otimes\hspace{-1pt}\hat{T}_{M}) \right) (I\hspace{-1pt}\otimes\hspace{-1pt} e_1\hspace{-5pt} \sqrt{2})\right\|\hspace{-2pt} 
&\leq 28 C_{M} \kappa(Z) M \left| \frac{c\gamma_\ell}{1-c} \right| ^{\ell + 1}, 
\end{align*}
concluding the proof.
\end{proof}

\end{document}
